\documentclass[11pt,letterpaper]{article}
\usepackage[T1]{fontenc}
\usepackage{lmodern,amsmath,amssymb,amsthm,microtype}
\usepackage[margin=1in]{geometry}
\usepackage[hidelinks]{hyperref}
\providecommand{\tightlist}{\setlength{\itemsep}{0pt}\setlength{\parskip}{0pt}}
\DeclareUnicodeCharacter{2264}{\ensuremath{\le}}
\DeclareUnicodeCharacter{2265}{\ensuremath{\ge}}
\DeclareUnicodeCharacter{2260}{\ensuremath{\ne}}
\DeclareUnicodeCharacter{2212}{\ensuremath{-}}

\DeclareUnicodeCharacter{2211}{\ensuremath{\sum}}

\usepackage{longtable,booktabs,array}
\title{Excluding eighth weights below 135 in the nine-weight search below 150}
\author{Aafi Khan Malek\\ \small Independent researcher, Ahmedabad, India\\
\small\href{mailto:aafimalek2032@gmail.com}{aafimalek2032@gmail.com}}
\date{28 November 2026}
\begin{document}
\maketitle
\textbf{Theorem.} Let nine sorted positive real weights have pairwise
unit-separated subset sums. If \(149\le x_9<150\), then \(x_8\ge135\).
The proof uses the preceding
\href{https://github.com/Aafimalek/novel-discoveries/tree/32b96614f57a26cbf30acb17235ad4baf40ff92a/nine-separated-subset-sums/verification/rank-cuts-below-150/penultimate-cap-133}{cap-133 theorem}
\(x_8\ge134\) and a complete exact replay of the next cap.
This does not determine the optimum of the nine-weight problem.

\section{Cap-specific implications}\label{cap-specific-implications}

The cap \(t=134\) forces coordinate-9 starts 145 through 172, so the
unsigned full-sign slots are \(q=37,\dots,64\). The magnitudes
\(\alpha_q=u_{a-q+1}\in[149-2q,151-2q)\) are consecutive, with common
\(64\le a\le74\) and \(b=148-a\). The established plus-two endpoints
then give the capacity ranges:

\begin{longtable}[]{@{}lr@{}}
\toprule\noalign{}
Half-open magnitude range & Capacity \\
\midrule\noalign{}
\endhead
\bottomrule\noalign{}
\endlastfoot
{[}1,21) & \(a-64\) \\
each {[}149−2q,151−2q), q from 64 down to 37 & 1 \\
{[}77,223) & 36 \\
{[}278,321) & \(148-2a\) \\
{[}375,10000) & \(a-20\) \\
\end{longtable}

The two wide gaps have width at least
\(\alpha_{37}-\alpha_{64}+4\ge58\). Flipping a sign at a weight strictly
below 29, or swapping signs at adjacent weights with a gap strictly
below 29, preserves membership at both cuts. The small-flip orbits on
complementary sign pairs give \(2^s\mid a,b\) for the number s of such
small weights, hence \(2^s\mid148\) and \(s\le2\). A nonsmall coordinate
exists because otherwise every absolute signed magnitude is below
\(8\cdot29<375\). Thus \(x_3\ge29\). Expanded zero-based positions
{[}0,a), {[}a,b), and {[}b,128) retain their exact capacities; an empty
middle group for a=b permits no middle label.

Let \(y=(128,\dots,135)\), \(r_i=y_i-x_i\), and \(R=\sum r_i\). Balanced
magnitudes at y are at most 16 and unbalanced magnitudes at least 248.
If some required form is unbalanced, then \(R>171\) because all required
magnitudes are below 77. Otherwise \(\alpha_{37}\ge75\) is balanced, so
\(R\ge59\). Therefore \(R\ge59\) in either case, and
\(r_3\le130-29=101\). The prior theorem supplies \(x_8\ge134\), hence
\textbf{\(r_8\le1\)}. For an ordered deficit tail of length N with sum
T, subtract \(r_8\) from each entry and use the extremal prefix-sign
averages of the first N−1 entries. The signed tail is enclosed by the
corresponding endpoint expressions \(\mu T+(\sum\sigma-N\mu)r_8\) over
the tail interval and \(0\le r_8\le\min(1,T_{\rm upper}/N)\), with
outward rounding.

\section{Deterministic coverage and exact
results}\label{deterministic-coverage-and-exact-results}

The 8,128 initial unit cells leave 7,223 before splitting by a. Of
79,453 two-deficit count profiles, 36,615 fail pure orbit size, 4,644
fail interval tests, and 21,070 fail the interval-conditioned orbit
assignment, leaving \textbf{17,124}. The exact canonical payload is
split into \textbf{140 disjoint shards} of at most 128 profiles; it
includes exactly every case that the first two stages did not reject.

All \textbf{140} clear with exact refinement: 102 at depth 3 and 38 at
depth 4. Summing the guarded worker logs gives:

\begin{center}
\small
\begin{tabular}{rrrrrr}
\toprule
Depth & Cells & Orbit size & Intervals & Orbit + intervals & Surviving \\
\midrule
3 & 825,854 & 199,181 & 110,392 & 510,489 & 5,792 \\
4 & 191,000 & 131,157 & 24,219 & 35,624 & 0 \\
\bottomrule
\end{tabular}
\end{center}

A lightweight independent audit checks the exact disjoint partition,
source pins, and the upstream and downstream digest of each completed
depth. The separate arithmetic audit reconciles all 178 recorded depth
logs against the exact progress counts. An independent fresh run can
recompute the actual proof arithmetic.

\section{Replay}\label{replay}

\begin{verbatim}
sha256sum -c SHA256SUMS
python3 parameters134.py
python3 coverage_audit.py
python3 arithmetic_totals.py
./replay-all.sh
\end{verbatim}

The final script copies pinned source and checkpoint to a fresh
temporary directory, reruns phase 2 and checks canonical byte equality,
and recomputes all 140 shards serially with no reuse of old progress.
Each proof invocation is bounded by 35 seconds and 143,000 KiB of
virtual memory; memory and disk reserves are checked between shards. All
the computations are standard Python integer or fraction arithmetic. The
geometric implications are handwritten; this is not a Lean theorem or an
exclusion of any candidate with \(x_8\ge135\).

\end{document}
